\documentclass{article}
\usepackage{graphicx}
\usepackage{float}
\usepackage{wrapfig}

% --- Encoding & Language ---
\usepackage[utf8]{inputenc}     % pdfLaTeX
\usepackage[T1]{fontenc}
\usepackage[magyar]{babel}

% --- Math & Chemistry ---
\usepackage{amsmath, amsthm, amssymb}
\usepackage[version=4]{mhchem} % chemical formulas
\usepackage{siunitx}            % units, temperatures
\usepackage{tikz}
\usetikzlibrary{tikzmark, calc}
\usepackage{mathtools}

% --- Section formatting ---
\usepackage{titlesec}

% Allow deep section numbering and TOC
\setcounter{secnumdepth}{5}
\setcounter{tocdepth}{5}

\makeatletter

% ------------------------
% LEVEL 4: subsubsubsection
% ------------------------
\titleclass{\subsubsubsection}{straight}[\subsubsection]
\newcounter{subsubsubsection}[subsubsection]
\renewcommand\thesubsubsubsection{\thesubsubsection.\arabic{subsubsubsection}}

\renewcommand\subsubsubsection{%
  \@startsection{subsubsubsection}{4}{\z@}%
  {3.25ex \@plus1ex \@minus.2ex}%
  {1.5ex \@plus .2ex}%
  {\normalfont\normalsize\bfseries}%
}

\titleformat{\subsubsubsection}
  {\normalfont\normalsize\bfseries}
  {\thesubsubsubsection}{1em}{}

\titlespacing*{\subsubsubsection}{0pt}{3ex}{1ex}

% TOC entry for level 4
\newcommand*\l@subsubsubsection{\@dottedtocline{4}{7em}{4em}}

% --------------------------
% LEVEL 5: subsubsubsubsection
% --------------------------
\titleclass{\subsubsubsubsection}{straight}[\subsubsubsection]
\newcounter{subsubsubsubsection}[subsubsubsection]
\renewcommand\thesubsubsubsubsection{\thesubsubsubsection.\arabic{subsubsubsubsection}}

\renewcommand\subsubsubsubsection{%
  \@startsection{subsubsubsubsection}{5}{\z@}%
  {3.25ex \@plus1ex \@minus.2ex}%
  {1.5ex \@plus .2ex}%
  {\normalfont\normalsize\bfseries}%
}

\titleformat{\subsubsubsubsection}
  {\normalfont\normalsize\bfseries}
  {\thesubsubsubsubsection}{1em}{}

\titlespacing*{\subsubsubsubsection}{0pt}{3ex}{1ex}

% TOC entry for level 5
% Use level 4 internally for safety (LaTeX only knows TOC levels 1-4)
\newcommand*\l@subsubsubsubsection{\@dottedtocline{4}{10em}{5em}}

\makeatother

\title{A fémek}
\author{Rovenszky Robin, Haydn Lorenzo}
\date{2026. Szeptember}

\newtheorem{definition}{Definíció}

\begin{document}

\maketitle
\tableofcontents

\vspace{14pt}
\hrule

\section{Az ötvözetek}

Például
\begin{itemize}
    \item Bronz $->$ réz $+$ ón
    \item Sárgaréz $->$ réz $+$ cink
    \item Acél $->$ vas $+$ szén
    \item \emph{Amalgám} $->$ higany $+$ egyéb fémek
    \item Alumínium ötvözetek (pl. Dúralumínium) $->$ Al; Mg; Li 
\end{itemize}

\noindent
Az ötvözetek előállításának célja: Új, jobb tulajdonságú anyagok létrehozása.

\begin{definition}
    Fémek, illetve nemfémek összeolvasztásával keletkező anyag.
\end{definition}

\section{Az alkálifémek}
I.A. Főcsoport, fémei:
\begin{itemize}
    \item Lítium (Li) $<-$ akkumulátorok
    \item Nátrium (Na) $<-$ fontos élettani hatás
    \item Kálium (K) $<-$ fontos élettani hatás
    \item Rubídium (Rb)
    \item Cézium (Cs)
    \item Francium (Fr)
\end{itemize}

\noindent
1 db vegyérték elektron $=>$ kation-képződés során leadják ($\text{Li}^+; \text{Na}^+; \text{K}^+$.
Nagyon reakcióképes elemek $->$ nem fordulnak elő elemi formában a természetben.

\subsection{Lítium}
Akkumulátorok, lítium bányászat, vagy folyókból kinyerik, elektromos autók

\subsection{Nátrium}
\subsubsection{Fizikai tulajdonságai}
\begin{itemize}
    \item Ezüst szürke fém, a felületét oxidréteg ($->$ petróleumban tároljuk) borítja
    \item Puha, késsel könnyen vágható
    \item Úszik a vízben (mivel $\rho(Na) < \rho(\text{víz})) $(és reagál is vele)
\end{itemize}

\subsubsection{Kémiai tulajdonságai}
\subsubsection{Kitekintő: Kísérlet}
\subsubsection{Előfordulása}
\subsubsection{Előállítása}
\subsubsection{Felhasználása}
\subsubsection{Élettani hatása}
\subsubsection{Vegyületei}

% Várjunk csak, most ez a káliumról is beszél (lsd. irl kézi jegyzet) vagy csak a nátriumról, és majd vesszük utánna a káliumot?

\end{document}